% ===== PRÉAMBULE CANONIQUE — à coller VERBATIM en tête de CHAQUE .tex =====
\documentclass[11pt,a4paper]{article}
\usepackage[utf8]{inputenc}\usepackage[T1]{fontenc}\usepackage{lmodern}
\usepackage[french]{babel}
\usepackage{amsmath,amssymb,mathtools}\usepackage{icomma}
\usepackage{siunitx}\sisetup{locale=FR}  % \qty{}{}, \si{}, \ang{}
\usepackage{xcolor}\usepackage{colortbl}
\usepackage{tikz}\usepackage{pgfplots}\pgfplotsset{compat=1.18}
\usepackage[siunitx,europeanresistors]{circuitikz} % circuits (élec.)
\usepackage[most]{tcolorbox}\usepackage{enumitem}\usepackage{array,booktabs}
\usepackage[a4paper,margin=2.2cm]{geometry}
\usetikzlibrary{arrows.meta,calc,angles,quotes,positioning,patterns,%
  backgrounds,shapes.geometric,decorations.pathmorphing,decorations.markings,intersections}
\definecolor{primary}{RGB}{222,49,99}\definecolor{primarydark}{RGB}{150,22,55}
\definecolor{defcol}{RGB}{0,114,178}\definecolor{methcol}{RGB}{52,92,148}
\definecolor{excol}{RGB}{0,135,90}\definecolor{attncol}{RGB}{200,40,55}
\tcbset{boxbase/.style={enhanced,breakable,boxrule=0.9pt,arc=2.5pt,
  left=3mm,right=3mm,top=2.3mm,bottom=2.3mm,fonttitle=\bfseries,coltitle=white}}
% --- boîtes TUTORIEL (noms EXACTS, ne pas renommer) ---
\newtcolorbox{cle}[1][]{boxbase,colback=defcol!6,colframe=defcol,
  title={\textsc{À retenir}\ifx\relax#1\relax\relax\else\textmd{ : \itshape #1}\fi}}
\newtcolorbox{methode}[1][]{boxbase,colback=methcol!7,colframe=methcol,sharp corners,
  title={\textsc{Méthode}\ifx\relax#1\relax\relax\else\textmd{ --- \itshape #1}\fi}}
\newtcolorbox{exemplebox}[1][]{boxbase,colback=excol!5,colframe=excol,
  title={\textsc{Exemple}\ifx\relax#1\relax\relax\else\textmd{ : \itshape #1}\fi}}
\newtcolorbox{piege}[1][]{boxbase,colback=attncol!6,colframe=attncol,
  title={\textsc{Piège}\ifx\relax#1\relax\relax\else\textmd{ : \itshape #1}\fi}}
\newtcolorbox{remarque}[1][]{boxbase,colback=black!4,colframe=black!45,
  title={\textsc{Remarque}\ifx\relax#1\relax\relax\else\textmd{ : \itshape #1}\fi}}
% --- boîtes EXERCICE / CORRIGÉ (noms EXACTS, auto-comptées) ---
\newtcolorbox[auto counter]{exo}[1][]{boxbase,colback=primary!4,colframe=primary,
  title={\textsc{Exercice}~\thetcbcounter\ifx\relax#1\relax\relax\else\textmd{ --- \itshape #1}\fi}}
\newtcolorbox[auto counter]{corrige}[1][]{boxbase,colback=excol!4,colframe=excol,
  title={\textsc{Corrigé}~\thetcbcounter\ifx\relax#1\relax\relax\else\textmd{ --- \itshape #1}\fi}}
\newcommand{\vv}[1]{\vec{#1}}\newcommand{\ds}{\displaystyle}
\newcommand{\R}{\mathbb{R}}\newcommand{\N}{\mathbb{N}}\newcommand{\Z}{\mathbb{Z}}
% (un \usepackage{fancyhdr}+en-tête est possible ; il est ignoré côté web)
% =========================================================================
\begin{document}

\begin{center}
{\large\bfseries\color{primarydark} Thème 2 — Niveau Moyen}\\[-2pt]
{\color{primary}\rule{5cm}{1pt}}
\end{center}

\begin{exo}[écrire l'équation d'une onde]
Une onde sinusoïdale d'amplitude $A=\SI{3}{\centi\metre}$ et de fréquence $f=\SI{25}{\hertz}$
se propage vers les $x>0$ à la célérité $v=\SI{10}{\metre\per\second}$.
\begin{enumerate}[label=\alph*)]
  \item Calcule la pulsation $\omega$.
  \item Calcule la longueur d'onde $\lambda$, puis le nombre d'onde $k$.
  \item Écris l'équation $y(x,t)$ de l'onde (en unités SI, phase à l'origine nulle).
\end{enumerate}
\end{exo}

\begin{exo}[tout tirer d'une équation]
Une onde a pour équation (unités SI) :
\[ y(x,t)=0{,}02\,\sin\!\left(100\pi\,t - 5\pi\,x\right). \]
\begin{enumerate}[label=\alph*)]
  \item Donne l'amplitude $A$, la pulsation $\omega$, puis la fréquence $f$ et la période $T$.
  \item Donne le nombre d'onde $k$, puis la longueur d'onde $\lambda$.
  \item Calcule la célérité $v$ de l'onde (deux méthodes possibles : $\lambda f$ ou $\omega/k$).
\end{enumerate}
\end{exo}

\begin{exo}[déphasage dans le temps]
Deux points d'une corde vibrent à la même période $T=\SI{2}{\second}$. Le second reproduit le
mouvement du premier avec un retard $\Delta t=\SI{0.5}{\second}$.
\begin{enumerate}[label=\alph*)]
  \item Calcule le déphasage $\Delta\varphi=2\pi\,\Delta t/T$ entre les deux points.
  \item Ces deux points sont-ils en concordance, en opposition, ou dans aucun de ces deux cas ?
\end{enumerate}
\end{exo}

\begin{exo}[déphasage dans l'espace]
Sur une même onde de longueur d'onde $\lambda=\SI{6}{\metre}$, on considère deux points séparés
d'une distance $d$. On rappelle $\Delta\varphi=2\pi\,d/\lambda$.
\begin{enumerate}[label=\alph*)]
  \item Pour $d=\SI{1.5}{\metre}$, calcule $\Delta\varphi$ et conclus sur la nature du déphasage.
  \item Pour $d=\SI{3}{\metre}$, calcule $\Delta\varphi$ et conclus.
\end{enumerate}
\end{exo}

\begin{exo}[le point va-t-il plus vite que l'onde ?]
Un point d'une corde oscille avec l'amplitude $A=\SI{5}{\centi\metre}$ et la fréquence
$f=\SI{20}{\hertz}$. L'onde se propage le long de la corde à la célérité $v=\SI{8}{\metre\per\second}$.
\begin{enumerate}[label=\alph*)]
  \item Calcule la vitesse maximale $v_{\max}$ du point du milieu.
  \item Compare $v_{\max}$ à la célérité $v$ de l'onde. Sont-ce les mêmes vitesses ?
\end{enumerate}
\end{exo}

\end{document}
